\section{Discussion}

This paper does not provide an empirical verdict and should not be considered a universal replacement. Its primary contribution is a decision framework for choosing educational AI architectures: classical RAG might still be the best top-level abstraction for simple, evidence-based tasks, while tasks requiring heavy coordination might benefit from agentic or multi-agent orchestration. In many mixed-regime environments, the most justifiable conclusion is often a hybrid approach rather than an absolute one, with retrieval as an available capability within a broader orchestration strategy.

This framing also emphasizes key limitations. Current benchmarks might not fully reflect the quality of genuine tutoring, learner modeling, or long-term instructional success. Multi-agent systems could lead to increased costs, more complex debugging, and new risks of coordination failure. Therefore, the paper focuses on architectural choices and evaluation design rather than making broad claims about classroom effects. While output quality, citation support, and adaptive behaviors can be informative, they should not be regarded as direct proof of learner improvements without additional research.

The main argument is therefore comparative rather than doctrinal. The key question isn’t just whether retrieval should be used, but whether a traditional retrieve-then-generate approach is enough for tasks requiring pedagogical sequencing, critique, adaptation, and role specialization. An effective evaluation should identify where classical RAG excels, where orchestration provides meaningful improvements, and where increased complexity mainly raises costs and latency without sufficient educational advantages. In this sense, the paper's contribution is not a final judgment but a systematic way to determine when retrieval-focused, orchestration-heavy, or hybrid systems are most appropriate.